\documentclass{article}

\usepackage[margin=1.25in]{geometry}

% Bibliography support
\usepackage[
	backend=biber,
	style=authoryear
]{biblatex}

% Load references from a separate .bib file
\addbibresource{Bibliography.bib}

% Custom section numbering
\renewcommand{\thesection}{(T\arabic{section})}
% Making it toc compatible
\makeatletter
\renewcommand{\l@section}{\@dottedtocline{1}{1.5em}{3.5em}}
\makeatother

%%% CUSTOM CITATION SETUP %%%

% Parenthetical citation: (P1)
\DeclareCiteCommand{\parencite}[\mkbibparens]
	{\usebibmacro{prenote}}
	{\printfield{entrykey}}
	{\multicitedelim}
	{\usebibmacro{postnote}}


% Text citation: Wohlin and Aurum (P1)
\DeclareCiteCommand{\textcite}
	{\boolfalse{citetracker}%
	\boolfalse{pagetracker}%
	\usebibmacro{prenote}}
	{\printnames{labelname}%
	\space\mkbibparens{\printfield{entrykey}}}
	{\multicitedelim}
	{\usebibmacro{postnote}}

% Sort bibliography by citation key
\DeclareSortingTemplate{byentrykey}{
	\sort{
		\field{entrykey}
	}
}


% Print citation key at the beginning of each entry
\renewbibmacro*{begentry}{%
	\printtext{\mkbibparens{\thefield{entrykey}}}%
	\hspace{1em}%
}

% Select the custom sorting template
\ExecuteBibliographyOptions{sorting=byentrykey}

%%%%%%%%%%%%%%%%%%%%%%%%%%%%%

\begin{document}

% Title page
\begin{titlepage}
    \centering
    \vspace*{4cm}
    {\LARGE\bfseries RMA - Research Methods Assignments\par}
    \vspace{1.3cm}
    {\large Master's Thesis in Technology Seminar, Software Engineering (DTEK1102)\par}
    \vspace{0.3cm}
    {\large Botond Ortutay\par}
    \vfill
    {\large \today\par}
\end{titlepage}

% Table of contents
\tableofcontents
\clearpage

\section*{Intro}
\addcontentsline{toc}{section}{Intro}

Your text goes here.

\begin{table}[htbp]
	\centering
	\begin{tabular}{|p{0.428in}|p{0.856in}|p{1.712in}|p{0.856in}|p{1.284in}|}
		\hline
		Code & Author & Title & University & Reason for choice \\
		\hline
		T1 & Lauri Kivimäki & Criteria for Sustainable Use of Artificial Intelligence & University of Turku & Deep learning topic \\
		\hline
		T2 & Olli Kiviniemi & A Method for Semi-Autonomous Map Generation and Path Finding In UAV-UGV Cooperation for Search and Rescue Mission in Dynamic and Unmapped Environment & University of Turku & Pathfinding topic \\
		\hline
		T3 & Juuso Haapanen & Katsaus transformer-arkkitehtuureihin ja niiden hyödyntämiseen konenäkötehtävissä & University of Turku & DL topic with CV applications \\
		\hline
		T4 & Sofia Sievinen & Requirements engineering practices for AI development – A case study in a mid-sized AI consulting company & Aalto University & Deep Learning in Software Engineering topic \\
		\hline
		T5 & Saku Lehtinen & Monitoring of real-time data pipelines & Tampere University & Data processing \& Software Engineering topic \\
		\hline
		T6 & Philippe de Bekker & AI for Software Engineering: Reviewing and Improving Benchmarking Practices & Delft University of Technology, Netherlands & AI for software engineering topic incl. benchmarking \\
		\hline
	\end{tabular}
	\caption{Assessment criteria}
	\label{tab:assessment-criteria}
\end{table}

\newpage
\section{Criteria for Sustainable Use of Artificial Intelligence}
% https://www.utupub.fi/server/api/core/bitstreams/cbb993ad-9b7c-4dd5-b8e1-13b887397569/content

% 6-8 bullet points total, at least one under each
% Research objective:
% - Finding out what kind of sustainability issues relate to Deep Learning & forming a set of actionable criteria to give to industry practicioners

\textcite{T1} explores the newly emerging Deep Learning (DL) technologies in 
his thesis from the perspective of how it could be made more sustainable in the 
future. He aims to first gather a realistic image of the current state of 
sustainability in the field and how they are taken into account in current 
practices of the industry. He then aims to form a set of actionable criteria 
professionals can use in the future to ensure their use of DL is as 
sustainable as possible.\parencite{T1}

% Type of Research (V2) according to *Evaluation Criteria/Red Flags*
% - Strategic research (subset of Basic): theoretical data collected to act as useful criteria in future practical use. Qualitative research: narrative text, observational data (as in: data observed in lit review), exporation of presented pre-existing research. Whenever numbers are presented they are there to support arguments rather than to serve a basis for quantitative analysis.

	\par I would classify the type of his research as strategic, which is 
a subset of basic research.\parencite{V2} This is because his way of 
conducting his research is mainly to gather as much theoretical data as he can 
from pre-existing sources, rather than conducting applied research in the 
traditional sense. I believe applied research into sustainable use of DL would 
look something like taking a typical DL problem, such as text or image 
generation for example and then building an application that is optimized for 
sustainability in some way. This process would then generate new knowledge on 
how to optimize DL systems for sustainability. However since the problem of 
\textcite{T1} is much more theoretical than this we are dealing with basic 
rather than applied research. But since the results are clearly aimed at 
practicioners to include their ways of working this research is strategic 
rather than pure basic. I would also classify this research as qualitative 
rather than quantitative or mixed, since it displays signs of quantitative 
research such as narrative text, an observational approach to the data 
presentation and an explorative approach to the gathering of the 
data.\parencite{V2} Whenever numbers are presented, they are there to support 
arguments rather than to serve a basis for quantitative analysis.

% Research questions:
% - RQ1: What is the current state of the field regarding the computational complexity, size and energy consumption of deep learning models and what are the trends of these metrics?
% - RQ2: Does the literature present estimates and measuring methods for the environmental impacts, energy consumption, or manufacturing carbon footprint of deep learning training and usage based on factors such as model size or architecture?
% - RQ3: Based on currently available research is it possible to develop big picture calculation and estimation criteria for optimizing deep learning for sustainability?

	\par The research questions of \textcite{T1} are clearly listed in 
chapter 1.1 of the thesis and are as follows:

\begin{description}
	\item[RQ1:] \label{t1rq1} What is the current state of the field regarding the computational complexity, size and energy consumption of deep learning models and what are the trends of these metrics?
	\item[RQ2:] \label{t1rq2} Does the literature present estimates and measuring methods for the environmental impacts, energy consumption, or manufacturing carbon footprint of deep learning training and usage based on factors such as model size or architecture?
	\item[RQ3:] \label{t1rq3} Based on currently available research is it possible to develop big picture calculation and estimation criteria for optimizing deep learning for sustainability?
\end{description}

% Research methods used:
% - Systematic literature review based on the PRISMA guideline; w/ published search string, inclusion/exclusion criteria; No dedicated methods section, lit review methodology in SLR Chapter intro (3, before 3.1); Possibly refer to (V4)
% - Data analysis: Chapter 4 attempts to form trend descriptions based on the knowledge collected in the SLR and to find problems to which chapter 5 attempts to find solution. However no explicit data analysis method has been named. Of the methods mentioned in (P1) thematic analysis is the closest since finding trends and problems can be thought of as coding the observed phenomena.
% - Artifact creation: Chapter 5's criteria for sustainable use of AI can be viewed as design science research (P1), however I don't see any observation being made on the actual application of the artifact which would be advisable in a design science project (V5). Then again this is Strategic research (V2) so maybe we shouldn't expect that either.

	\par To answer these questions several research methods have been 
employed. The most important of these, on top of which everything else is built 
is the Systematic Literature Review (SLR) present in chapter 3. Afterwards some 
basic data analysis is performed on the SLR results in chapter 4. Finally a set 
of 15 criteria for sustainable use of DL is presented in chapter 5.
	\par The SLR is based on the so-called PRISMA guideline,\parencite{E1} 
which is a guideline designed for better reporting of systematic reviews. By 
quickly glancing at the 2020 version of the PRISMA checklist, the work done by 
\textcite{T1} appears compliant to the PRISMA criteria. Additionally search 
strings and inclusion and exclusion criteria were published and I believe the 
methodology on how the SLR was conducted were described to a satisfactory 
degree. However I must note that this thesis doesn't have a dedicated 
methodology section and the relevant information was scattered throughout the 
different sections of the thesis.
	\par Chapter 4 attempts to form trend descriptions based on the 
knowledge collected in the SLR and to find problems to which chapter 5 attempts 
to find solutions. However no explicit data analysis method has been named. The 
data analysis however is undouptedly qualitative in it's nature and of the 
methods named by \textcite{P1} thematic analysis is the closest, since finding 
trends and problems can be thought of as coding the observed phenomena.
	\par Furthermore the creation of the set of sustainability criteria for 
DL model use can be thought of as artifact creation, and thus this work can be 
thought of as an example of design science.\parencite{P1} On the other hand 
the efect of the created artifact on the sustainability of DL usage is not 
then observed, would be advisable in a design science project.\parencite{V5} 
But then again we shouldn't expect that kind of approach from strategic 
research.\parencite{V2}

% Should some other RM have been used:
% - A named data analysis method would be nice. A lot of things are based on "the judgment of the author" in this thesis

	\par Over all, when it comes to chosen research methods, I believe that 
the SLR was a good choice and has been conducted in a satisfactory way, with a 
good enough reporting on the methodology. However I'd like to have seen 
additional data collection methods as well. All knowledge handled in this 
thesis came from pre-existing science, so no new measurements or data 
collection was involved in this thesis at all. Furthermore a named data 
analysis method would have been nice. The analysis itself as well as the 
criteria generation part is largely based on the judgment, so opinions, of the 
author. Although this bias in explicitly noted in the thesis itself.\parencite{T1}
	\par All in all the work of \textcite{T1} is basically a large and 
comprehensive literature review on what the Computer Science field says on the 
topic of sustainable usage of DL systems. On top of this there is the set of 
criteria for sustainable use of DL, which is again, broad and rooted in the 
work done on this field. However the quality of the analysis and specifically 
the biasedness of it could be improved.

\newpage
\section{A Method for Semi-Autonomous Map Generation and Path Finding In UAV-UGV Cooperation for Search and Rescue Mission in Dynamic and Unmapped Environment}
% https://www.utupub.fi/server/api/core/bitstreams/193b914f-34d4-4137-a681-13a67d67896d/content

% Main source for analysis: P2
% 6-8 bullet points total, (RQs count as one) at least one under each
% Research objective:
% - Allow UAV-UGV co-operarion for path finding in a search and rescue operation

In his thesis \textcite{T2} talks about a multi-agentic way to utilize both 
Unmanned Aerial Vehicles (UAV) and Unmanned Ground Vehicles (UGV) together in 
a search and rescue setting. More specifically he assumes a setup where one or 
more UGVs are sent on a search and rescue mission to a known end location in an 
unknown environment. The assumption is that search and rescue missions take 
place in disaster areas or war zones and thus even though it is known to which 
location the UGVs need to go, there is no reliable mapping data of the areas 
they'd need to traverse due to the dynamicity of the area. Therefore one or 
more UAVs are sent out to perform mapping to enable the UGVs to reach the 
target area. The thesis aims to examine the technical difficulties and 
feasibility of such a setup

% Research questions:
% - No traditional RQ or methodology section in thesis structure. No explicitly named RQs either. Even though 2.1 intro directly hints at existing RQs??? This is obviously an artifact creation first approach

	\par Analyzing the research methods used in the thesis of \textcite{T2} 
was more difficult than it should've been, because of its structure. There is 
no methodology section nor explicitly named research questions anywhere in this 
thesis. The only reference to research questions I was able to find in this 
thesis was a justification of the background and literature review portion of 
the text, where it claimed that a background chapter on each key technology 
discussed was necessary to understand their relation to each other and the 
research question.\parencite{T2} However still, research questions 
notwithstanding this is clearly an artifact creation first approach, since it 
is mentioned several times that the goal is to develop a new more robust 
mapping technology for the studied situation.\parencite{T2}

% - RQs could be something like:
%   - Given a dynamic environment an depth data captured from above:
%     1. How can we find a high quality (slopes, manuverability) path for an UGV so that it is able to operate?
%     2. How can we filter out false obstacles that will be present in the depth data? (foliage, powerlines etc)
%     3. How do we stitch together all various pieces of data generated by the UAV to form a coherent and realistic interpretation of the operating environment?

	\par Even though there are no explicitly defined research questions, 
by looking at what \textcite{T2} describe at the conclusion section, and the 
work that's actually described throughout the thesis, the research questions 
being answered in this work could be something like this:

\begin{itemize}
	\item Given a dynamic environment of which depth data is captured from above using an UAV:
	\begin{enumerate}
		\item How can we find a high quality path for an UGV so that it is able to operate?
		\item How can we filter out false obstacles that will be present in the depth data?
		\item How do we stitch together all various pieces of data generated by the UAV to form a coherent and realistic interpretation of the operating environment?
	\end{enumerate}
\end{itemize}

% Deriving suitable research methods from RQs based on Järvinen's taxonomy (P2):
%
% We are studying reality rather than proving a mathematical theorem.
% We are examining artifact utility rather than observing the reality of a situation
% We are building a new artifact to deal with an identified problem rather than evaluating an existing artifact (although there is a literature review section attempting to describe earlier work and they do note shortcomings & considerations from earlier work; still the main focus here is undouptadly artifact building rather than evaluation.
% A typical RQ for this kind of work would look like "Can a certain artifact be built?" and indeed all of our three hypothetical RQs follow this format (can a HQ path for an UGV be found? Can false obstacles be filtered out? Can all the data be stiched together?)
%

	\par Now let's use the taxonomy developed by \textcite{P2} backwards 
to see whether the hypothetical research questions presented here match the 
work conducted by \textcite{T2}. \textcite{T2} studies reality rather than 
trying to prove a mathematical theorem, so we're not dealing with a 
mathematical approach. \textcite{T2} deals with utilizing artifacts rather than 
observing the reality of a situation, so we're not dealing with an approach 
stressing what is reality. \textcite{T2} focuses on artifact-building more 
than artifact evaluation. There is a literature review section that does 
describe earlier similar solutions including their shortcomings from an 
evaluative perspective,\parencite{T2} but a much larger part of the thesis is 
focused on artifact-building and the various problems related to that process 
so that is the main focus approach of the thesis. \textcite{P2} says that 
research questions with the format \textit{"Can a certain artifact be built?"} 
should be mapped on to artifact-building approaches, which means that our 
hypothetical research questions fit the research well, which means that this 
work is indeed artifact-building-focused in its nature.

% Research methods used:
% - Artifact building (P2 analysis); Where the methodology/logic flow/pipeline itself is the artifact; The artifact is itself seen as a result (T2)

	\par The actual artifact built and presented in this thesis is quite 
abstract. \textcite{T2} himself describes it as 
\textit{a method for map generation and path finding in a dynamic and unmapped environment.}
In actuality this is a collection of pseudocode descriptions of several 
different algorithms that solve various parts of this huge overlapping 
problems as well as a detailed description of how these standalone algorithms 
should be made to interact with one another to provide a solution. This 
artifact, so the algorithms and their interaction as a flow, is then seen as 
the \textit{results} of his work.\parencite{T2}

% - Design science (V5); There is a clear and well thought out plan based on which the artifact is built

	\par One could also categorize this work as design science, although 
that term is not used anywhere within it. Still: in this work a problem is 
identified, a potential solution is planned by studying various aspects of the 
problem and mapping the solutions to a planned artifact, which is then 
implemented according to said plan, while the whole process is documented with 
scientific accuracy. This satisfies the criteria of design 
science.\parencite{V5}

% - Simulation used in actually testing out an implementation of all the different pieces of the artifact. Galliers and Land (E2) categorize simulation as a suitable research method for research where the object is technology.

	\par Although the artifact they present in this thesis is just 
algorithms and pseudocode there was an existing implementation also, for a 
simulated case where the algorithms in question were used is described in the 
thesis.\parencite{T2} This is treated as a proof of concept that the different 
algorithms described in the thesis can work together and output the desired 
type of data (so a map and a path through it). There are no quantitative 
measurements and studies performed on the simulation output. But examples of 
output are indeed provided and described to a satisfactory degree. While only 
testing this artifact in a simulated environment programmed in a game engine 
isn't a perfect portrayal of how it would perform in a real search and rescue 
mission, \textcite{E2} do categorize simulation as a suitable research mode for 
scientific research where the research object is technology, so this is a 
sensible choice by that criteria.

% - Literature review. Not mentioning framework, search strings, databases, article amounts, bases of source selection etc. Seemingly chose specific cases where a similar application was built. So the approach is similar to what's in my own thesis. Still not enough detail is given on the lit review.

	\par Finally I also have to mention that there is a literature review 
conducted as a part of this work. However it is rather subpar. There is no 
framework mentioned whose rules it is built on, no search strings or searched 
databases have been published, the question 
\textit{"How many articles were read as a part of the literature review?"} is 
unanswered and even the basis of selection of the chosen articles is not 
mentioned. The logic here seems to be that since this is an artifact-building 
based research project, the literature review should find as many similar 
projects as possible and learn from their approaches: How do they aim to 
tackle their problems at hand, what difficulties they run into, are there any 
major shortcomings in the reported results, etc. This approach is actually 
quite similar to my own thesis' literature review. Still, not enough details 
are given on the logic behind the literature review and it seems separate from 
the rest of this work.

% Should some other RM have been used:
% - Some in vitro simulation rather than all in silico (nothing runs on real hardware, major question marks remain (where to do processing)). Shortcomings acknowledged in thesis. 
% - Systematic lit review rather than unsystematic/feels like an afterthought

	\par In my opinion most of the issues of this work is that this is 
clearly a development project first and a science project second. There are no 
clear research questions and not much thought has been given to the 
methodology, but most of the thesis is spent describing the artifact, the 
development of it and the related technologies. But with that being said, when 
looking at this like a development project the selected methods 
(artifact-building, design science, simulation, implementation focused 
literature review) make sense. I would choose the same methods and I didn't 
really find any basis for different choices in the provided methodological 
sources either. That being said I do have issues with the certain ways this 
work has been carried out.
	\par To loan a term from biology: I don't like how all the simulations 
done for this work are \textit{in silico}. I believe it would've been feasible 
to set up some kind of \textit{in vitro} simulations. That is: Use real 
hardware and real objects (for example some building blocks in a room) instead 
of only building a fictional environment in a game engine, since now it is 
entirely possible that the approach doesn't take into account some factor that 
is obvious in a real environment with real hardware but is completely missed 
due to the simulation choices. On the other hand I have to note that these 
shortcomings were also acknowledged in the thesis as well, so the author is 
aware of this limitation.\parencite{T2} Another thing I would change is the 
nature of the literature review. I'd make it more systematic, based on some 
framework, report on it more thoroughly and integrate it better with the rest 
of the work. 
	\par To summarize: The work done by \textcite{T2} is a practical 
solution to a large problem, of which his work provides a large overview. I 
like the technical part of his work very much, but his work would benefit from 
better reporting and structure, and focusing on the scientific parts a bit 
more. The research method choices were sensible, but execution was imperfect. 

\newpage
\section{Katsaus transformer-arkkitehtuureihin ja niiden hyödyntämiseen konenäkötehtävissä}
% https://www.utupub.fi/server/api/core/bitstreams/5fa20755-a6fc-46e4-953e-c6e02b8b80ef/content

% Main source for analysis: V3
% 6-8 bullet points total, (RQs count as one), at least one under each
% Research objective:
% Research objective:
% - Given that:
%   - Computer vision is a disciplinary field, focusing on the interpretation of images and visual data using computer sysems
%   - Currently the best performing solutions for certain CV tasks are DCNNs
%   - However transformer models are used more and more to perofrm these types of tasks
%
% This thesis takes several transformer models developed for a CV task and compares them to each other as well as to pre-estabished DCNNs and explain methodoligacal and performance differences between them.

%\begin{displayquote}
%Konenäkö ja kuvien tunnistaminen tietokoneen avulla on erittäin pitkään tutkittu ongelma.
%\end{displayquote} (T3)

%\begin{displayquote}
%Kuvien tunnistaminen syvien neuroverkkojen avulla on koneoppimistehtävä, jossa parhaiten suorituvat menetelmät perustuvat tällä hetkellä konvoluutioneuroverkkoihin.
%\end{displayquote} (T3)

%\begin{displayquote}
%Transformer-arkkitehtuurit ovat alunperin luonnollisen kielen käsittelyyn suunniteltuja neuroverkkoja, mutta viimeaikainen tutkimus on osoittanut että niitä voidaan hyödyntää myös konenäön tehtäviin.
%\end{displayquote} (T3)

%\begin{displayquote}
%Tässä työssä tarkastellaan transformer-arkkitehtuureja, erityisesti konenäön tarpeisiin kehitettyjä transformer-malleja, sekä analysoidaan niiden menetelmällisiä ja suorituskyvyn eroja.
%\end{displayquote} (T3)
%%%%%%%%%%%%%%%%%%%%%%%%

% Research questions:

%\begin{displayquote}
%\begin{itemize}
%	\item TK1: Miten transformer-arkkitehtuuri eroaa konvoluutioneuroverkon arkkitehtuurista?
%\end{displayquote} (T3)

%\begin{displayquote}
%	\item TK2: Mikä on transformer-arkkitehtuurin pohjalta tehdyn neuroverkon suorituskyky verrattuna \textit{state-of-art} konvoluutioverkkoon kuvien luokittelussa ja semanttisessa segmentoinnissa?
%\end{itemize}
%\end{displayquote} (T3)


% Research methods used:
% - Both a literature review and an independent performance experiment (benchmark) is conducted.
% - The thesis claims to be an empirical study and a "theoretical observation/review" (tarkastelu). However an empirical study should either test or create a theory (P2) and this study doesn't clearly do either; RQs don't clearly state a theory they test. Nor are the RQs clearly theory creating ones either. RQ1/TK1 could be percieved as such but it doesn't really generate new information as such; It's more of a literature review kind of question; better fits conceptual-analytical from (P2). RQ2/TK2 does generate new knowledge and could be percieved as a theory creating question.
% - The literature review is used to answer RQ1/TK1, choose the appropriate neural networks for RQ2/TK2 and gather reported performance metrics for them. Search strings and databases were repoerted, however a more comprehensive lit review framework is not mentioned.

%\begin{displayquote}
%Transformer-arkkitehtuureita vertaillaan työssä konvoluutioneuroverkkoihin sekä kirjallisuudessa raportoitujen tuloksien perusteella sekä kokeellisin menetelmin.
%\end{displayquote} (T3)

%\begin{displayquote}
%pitkälti empiiristä tutkimusta sekä teoreettista tarkastelua. Empiirinen tukimusmenetelmä perustuu tutkimuskohteen havainnointiin ja mittaamiseen, joka tässä työssä tarkoittaa tutkittavien algoritmien hyvyyden arviointia tietyn aineiston ja ennustamistehtävän ratkaisussa.
%\end{displayquote} (T3)

%\begin{displayquote}
%When the past and present are empirically studied, we differentiate the theory-testing or theory-creating methods depending on whether there is a theory, model or framework guiding the study or is a researcher developing a new theory grounded on the gathered raw data.
%\end{displayquote} (P2)

%\begin{displayquote}
%Tähän työhön valittiin eri arkkitehtuureita tekemällä haku Scopus-tietokantaan \textit{"image classification"AND transformers.}
%\end{displayquote} (T3)

% - The empirical experiment (benchmark) in this thesis was: two pre-trained transformers (ViT, Swin Transformer) and a pre-trained DCNN (EfficientNet B0), were fine-tuned on an image classification task on the Cifar100 dataset. Results are reported as both raw numbers (accuracy, percentage) and as a written narrative. This is a sign of the usage of mixed (both qualitative and quantitative) research methods.

%\begin{displayquote}
%Valitut mallit olivat \textsc{ViT, Swin Transformer} sekä \textsc{EfficientNet B0}, jotka ovat kirjallisuudessa esitellyiltä tuloksiltaan hyviä. Vertailuun valittu ongelma oli kuvien luokittelu eri luokkiin. Kuvien luokittelu valittiin vertailun tehtäväksi siksi, että se on konenäön tehtävistä yksinkertaisin. Aineistona testissä käytetään Cifar100-aineistoa
%\end{displayquote} (T3)

%\begin{displayquote}
%Kokeelliset tulokset osoittavat, että transformer-pohjaiset mallit ovat kilpailukykyisiä verrattuna konvoluutioneuroverkkoihin. Tässä työssä hienosäädetty ViT-malli saavutti 91,2\% ja Swin Transformer 85,78\% tarkkuuden, kun taas EfficientNet-malli, joka edustaa konvoluutioneuroverkkoja, saavutti 82,01\% tarkkuuden kuvien luokittelutehtävässä.
%\end{displayquote} (T3)

%\begin{displayquote}
%Aluksi otetaan käyttöön esikoulutettu malli, jossa on valmiiksi laskettuna painot tietylle aineistolle, tässä tapauksessa ImageNet-aineistolle. Muut kuin mallin päällikerros jäähdytetään (engl. \textit{freeze}) siten, että sen painot eivät voi muuttua koulutuksen aikana vaan ainoastaan päällimmäisen kerroksen painot lasketaan uudestaan. Näin saadaan laskettua ennustetukset. Päällikerroksen kouluttamisen jälkeen koko malli voidaan kouluttaa uudelleen, mutta tässä työssä tämä jätettiin tekemättä. Valmiina malleina käytettiin kirjastoista saatavia valmiita malleja.
%\end{displayquote} (T3)

%\begin{displayquote}
%Mallien hienosäädön perusteella, ilman kattavaa hyperparametrien säätämistä voidaan todeta että sekä \textsc{ViT} että \textsc{Swin Transformer}-pohjaisen mallin tarkkuus on parempi kuin konvoluutioverkkoihin perustuvan mallin. Kirjallisuudessa esitettyjen tuloksien, jotka on esitetty taulukossa 5.2, perusteella EfficientNet-mallin pitäisi päihittää Swin Transformer, koska käytössä oli mallin ensimmäinen versio. Tässä kokeellisessa osuudessa ei kuitenkaan ole käytetty samoja parametreja kuin kirjallisuudessa, joten eroja teoreettisiin tuloksiin voi esiintyä. Yleisesti ottaen sekä kokeellisten että kirjallisuudessa esiintyvien tuloksien perusteella voi arvioida, että sekä transformer- että konvoluutioneuroverkot pärjäävät hyvin luokittelutehtävässä. Molempien mallien koulutuksissa merkille pantavaa oli mallin kouluttamisen kesto, joka oli yli 20 minuuttia kaikilla malleilla.
%\end{displayquote} (T3)

%\begin{displayquote}
%mixed method may involve analyzing both numerical and non-numerical data it may also present results in tables and graphs backed up by a strong narrative
%\end{displayquote} (V2)

%%% V3 %%%
% Examples of event/phenomena descriptions/interpretations in the text
% Any methods mentioned in V3 used?
% Steps of qualitative analysis described in V3 notes present?
% V3/4 quality analysis
%%% %% %%%

% Should some other RM have been used:

\newpage
\section{Requirements engineering practices for AI development – A case study in a mid-sized AI consulting company}
% https://aaltodoc.aalto.fi/server/api/core/bitstreams/f18a1d5d-1b24-414b-853b-7111caf85eff/content

% Main source for analysis: V1
% 6-8 bullet points total, (RQs count as one), at least one under each
% Research objective:
% Research questions:
% Research methods used:
% Should some other RM have been used:

\newpage
\section{Monitoring of real-time data pipelines}
% https://trepo.tuni.fi/bitstream/handle/10024/233097/LehtinenSaku.pdf

% Main source for analysis: V4
% 6-8 bullet points total, (RQs count as one), at least one under each
% Research objective:
% Research questions:
% Research methods used:
% Should some other RM have been used:

\newpage
\section{AI for Software Engineering: Reviewing and Improving Benchmarking Practices}
% https://repository.tudelft.nl/record/uuid:826586d1-b588-4a4e-9fc9-bb3d62f521ce
% Broken link. Point to git repo if no working copy found

% Main source for analysis: V3
% 6-8 bullet points total, (RQs count as one), at least one under each
% Research objective:
% Research questions:
% Research methods used:
% Should some other RM have been used:

\newpage
\section*{My thesis}
\addcontentsline{toc}{section}{My thesis}

Your text goes here.



% Print the bibliography
\newpage
\printbibliography[heading=bibintoc,title={References}]







\newpage
\appendix

\section{On use of generative AI}

GPT + Gemini: Notes on reference material. 
GPT: as a search engine for finding relevant MSc Theses from non-Finnish 
universities (for Finnish universities UTUPub + Trepo + Aaltodoc + Lutpub 
searched directly, no AI used in finding theses)
GPT: Latex environment setup. 
No genAI used in: thesis analysis, report skeleton generation, text generation

\end{document}
